% ============================================================
% Close-up: BioMarkerKGAgent — Workflow Interface View
% Fragment — % ============================================================
% Close-up: BioMarkerKGAgent — Workflow Interface View
% Fragment — % ============================================================
% Close-up: BioMarkerKGAgent — Workflow Interface View
% Fragment — % ============================================================
% Close-up: BioMarkerKGAgent — Workflow Interface View
% Fragment — \input{figs/closeup_biomarkerkg}
% ============================================================
\begin{figure}[htbp]
\centering
\resizebox{\textwidth}{!}{%
\begin{tikzpicture}[node distance=0.35cm, font=\sffamily,
    stepbar/.style={rectangle, rounded corners=6pt, thick,
        minimum width=19cm, align=center, font=\sffamily,
        drop shadow={opacity=0.10}, inner sep=6pt, text width=18cm},
    rwpanel/.style={rectangle, rounded corners=5pt, thick,
        minimum width=19cm, align=left, inner sep=6pt,
        font=\sffamily\small, text width=17.5cm},
    infopanel/.style={rectangle, rounded corners=5pt, draw=gray!50, thick,
        minimum width=19cm, align=left, fill=white,
        inner sep=6pt, font=\sffamily\small, text width=17.5cm},
    toolpanel/.style={rectangle, rounded corners=5pt, draw=gray!60, thick,
        minimum width=19cm, align=left, fill=gray!4,
        inner sep=6pt, font=\sffamily\small, text width=17.5cm},
    flowline/.style={->, >=Stealth, very thick, draw=gray!40},
    consumerlabel/.style={font=\sffamily\small\itshape, text=gray!60!black,
        align=center, text width=19cm}
]
    % === STEP BAR ===
    \node[stepbar, draw=cyan!50!black, fill=agentCyan] (step)
        {\iconGraph{cyan!60!black}\;\;
         \textbf{\large BioMarkerKGAgent --- Knowledge-Graph
          Contextualisation}\\[3pt]
         {\small Planner: \textit{``Query top 10--15 DEGs for KG
          neighbours --- builds Gene$\rightarrow$Pathway$\rightarrow$Phenotype
          chains.''}}\\[2pt]
         {\scriptsize\sffamily\color{gray!60!black}
          Shown as a second step in this example;
          the planner may invoke it earlier, later, or
          re-invoke it after replanning.}};

    % === READS ===
    \node[rwpanel, draw=green!40!black, fill=agentGreen,
          below=0.4cm of step] (reads)
        {\textbf{$\blacktriangledown$ Reads from AgentState}\\[5pt]
         $\bullet$\; \texttt{query} {\scriptsize (str)} ---
          original research question\\[2pt]
         $\bullet$\; \texttt{top\_genes} {\scriptsize (list[str])} ---
          ranked DEGs \emph{from OmicMiningAgent}\\[2pt]
         $\bullet$\; \texttt{previous\_results} {\scriptsize (str)} ---
          accumulated text output from Step~1};

    \draw[flowline] (step.south) -- (reads.north);

    % === PROMPT HIGHLIGHTS ===
    \node[infopanel, below=0.4cm of reads] (prompt)
        {\textbf{System Prompt Highlights}
         {\scriptsize (from \texttt{BioMarkerKGAgent.system\_message})}\\[5pt]
         \textbf{1.}\; ``Query the PrimeKG knowledge graph to retrieve
          relevant biomedical information''\\[2pt]
         \textbf{2.}\; ``Provide structured data and relationships from
          the knowledge graph to enhance understanding''\\[2pt]
         \textbf{3.}\; {\itshape (From planning prompt)}\;
          Query for \emph{both} up-regulated and down-regulated gene
          sets\\[2pt]
         \textbf{4.}\; {\itshape (From planning prompt)}\;
          Be THOROUGH --- query KG for multiple gene sets to build
          Gene$\rightarrow$Pathway$\rightarrow$Phenotype chains};

    \draw[flowline] (reads.south) -- (prompt.north);

    % === TOOL ===
    \node[toolpanel, below=0.4cm of prompt] (tool)
        {\textbf{Tool:}\;
         \texttt{biomarker\_kg\_tool(query)}\\[5pt]
         {\scriptsize \textbf{Internal pipeline:}\;
          text embedding $\rightarrow$
          PrimeKG subgraph retrieval (PCST) $\rightarrow$
          GAT graph encoder $\rightarrow$
          Llama-3.1-8B decoder $\rightarrow$
          ranked hypotheses}\\[3pt]
         {\scriptsize \textbf{Returns:}\;
          drug interactions, pathway memberships, GO terms,
          protein--protein interactors, disease associations}\\[3pt]
         {\scriptsize \textbf{Resources:}\;
          Neo4j/PrimeKG (20k+ diseases, 30k+ proteins)\; $\cdot$\;
          GRetriever micro-service (port~8001)\; $\cdot$\;
          OpenAI embeddings\; $\cdot$\;
          PyTorch Geometric}\\[3pt]
         {\scriptsize \textit{See companion illustration for internal
          GRetriever architecture.}}};

    \draw[flowline] (prompt.south) -- (tool.north);

    % === WRITES ===
    \node[rwpanel, draw=blue!30!black, fill=stateBlue,
          below=0.4cm of tool] (writes)
        {\textbf{$\blacktriangle$ Writes to AgentState}\\[5pt]
         $\bullet$\; \texttt{agent\_outputs[]} {\scriptsize (str)} ---
          structured KG results per gene, including:\\[2pt]
         \hphantom{$\bullet$\;}
          {\scriptsize --- pathway membership chains
          (e.g.\ \textit{CDK5 $\rightarrow$ cell-cycle signalling
          $\rightarrow$ neuronal death})}\\[2pt]
         \hphantom{$\bullet$\;}
          {\scriptsize --- GO biological-process terms}\\[2pt]
         \hphantom{$\bullet$\;}
          {\scriptsize --- drug--gene interactions and existing
          therapeutic targets}\\[2pt]
         \hphantom{$\bullet$\;}
          {\scriptsize --- protein--protein interaction partners}};

    \draw[flowline] (tool.south) -- (writes.north);

    % === DOWNSTREAM CONSUMERS ===
    \node[consumerlabel, below=0.3cm of writes] (consumers)
        {$\downarrow$\; Consumed by:\;
         \textbf{PubMedResearcher} (pathway context)\; $\cdot$\;
         \textbf{ScientistsAgent} (KG evidence for hypothesis
          anchoring)\; $\cdot$\;
         \textbf{ReportingNode} (KG analysis section)};

\end{tikzpicture}%
}% end resizebox
\caption{\textbf{BioMarkerKGAgent} --- workflow interface view.
In the typical realisation shown here this agent runs second, but
the LLM planner can reorder, skip, or re-invoke any agent at any
point---for instance, the planner may call it again after literature
search reveals new gene targets.  When invoked, it reads the
\texttt{top\_genes} list from \texttt{AgentState} and queries the
PrimeKG knowledge graph for each gene's neighbourhood, targeting the
top 10--15 DEGs to build explicit
Gene$\rightarrow$Pathway$\rightarrow$Phenotype chains.  Internally,
the tool communicates with a GRetriever micro-service that extracts
PCST subgraphs, encodes them with a Graph Attention Network, and
decodes answers with Llama-3.1-8B (see companion illustration for
architecture details).  Its output---appended to
\texttt{agent\_outputs[]}---provides the pathway and interaction
context that downstream agents use to anchor hypotheses in structured
biomedical knowledge.}
\label{fig:closeup-biomarkerkg}
\end{figure}

% ============================================================
\begin{figure}[htbp]
\centering
\resizebox{\textwidth}{!}{%
\begin{tikzpicture}[node distance=0.35cm, font=\sffamily,
    stepbar/.style={rectangle, rounded corners=6pt, thick,
        minimum width=19cm, align=center, font=\sffamily,
        drop shadow={opacity=0.10}, inner sep=6pt, text width=18cm},
    rwpanel/.style={rectangle, rounded corners=5pt, thick,
        minimum width=19cm, align=left, inner sep=6pt,
        font=\sffamily\small, text width=17.5cm},
    infopanel/.style={rectangle, rounded corners=5pt, draw=gray!50, thick,
        minimum width=19cm, align=left, fill=white,
        inner sep=6pt, font=\sffamily\small, text width=17.5cm},
    toolpanel/.style={rectangle, rounded corners=5pt, draw=gray!60, thick,
        minimum width=19cm, align=left, fill=gray!4,
        inner sep=6pt, font=\sffamily\small, text width=17.5cm},
    flowline/.style={->, >=Stealth, very thick, draw=gray!40},
    consumerlabel/.style={font=\sffamily\small\itshape, text=gray!60!black,
        align=center, text width=19cm}
]
    % === STEP BAR ===
    \node[stepbar, draw=cyan!50!black, fill=agentCyan] (step)
        {\iconGraph{cyan!60!black}\;\;
         \textbf{\large BioMarkerKGAgent --- Knowledge-Graph
          Contextualisation}\\[3pt]
         {\small Planner: \textit{``Query top 10--15 DEGs for KG
          neighbours --- builds Gene$\rightarrow$Pathway$\rightarrow$Phenotype
          chains.''}}\\[2pt]
         {\scriptsize\sffamily\color{gray!60!black}
          Shown as a second step in this example;
          the planner may invoke it earlier, later, or
          re-invoke it after replanning.}};

    % === READS ===
    \node[rwpanel, draw=green!40!black, fill=agentGreen,
          below=0.4cm of step] (reads)
        {\textbf{$\blacktriangledown$ Reads from AgentState}\\[5pt]
         $\bullet$\; \texttt{query} {\scriptsize (str)} ---
          original research question\\[2pt]
         $\bullet$\; \texttt{top\_genes} {\scriptsize (list[str])} ---
          ranked DEGs \emph{from OmicMiningAgent}\\[2pt]
         $\bullet$\; \texttt{previous\_results} {\scriptsize (str)} ---
          accumulated text output from Step~1};

    \draw[flowline] (step.south) -- (reads.north);

    % === PROMPT HIGHLIGHTS ===
    \node[infopanel, below=0.4cm of reads] (prompt)
        {\textbf{System Prompt Highlights}
         {\scriptsize (from \texttt{BioMarkerKGAgent.system\_message})}\\[5pt]
         \textbf{1.}\; ``Query the PrimeKG knowledge graph to retrieve
          relevant biomedical information''\\[2pt]
         \textbf{2.}\; ``Provide structured data and relationships from
          the knowledge graph to enhance understanding''\\[2pt]
         \textbf{3.}\; {\itshape (From planning prompt)}\;
          Query for \emph{both} up-regulated and down-regulated gene
          sets\\[2pt]
         \textbf{4.}\; {\itshape (From planning prompt)}\;
          Be THOROUGH --- query KG for multiple gene sets to build
          Gene$\rightarrow$Pathway$\rightarrow$Phenotype chains};

    \draw[flowline] (reads.south) -- (prompt.north);

    % === TOOL ===
    \node[toolpanel, below=0.4cm of prompt] (tool)
        {\textbf{Tool:}\;
         \texttt{biomarker\_kg\_tool(query)}\\[5pt]
         {\scriptsize \textbf{Internal pipeline:}\;
          text embedding $\rightarrow$
          PrimeKG subgraph retrieval (PCST) $\rightarrow$
          GAT graph encoder $\rightarrow$
          Llama-3.1-8B decoder $\rightarrow$
          ranked hypotheses}\\[3pt]
         {\scriptsize \textbf{Returns:}\;
          drug interactions, pathway memberships, GO terms,
          protein--protein interactors, disease associations}\\[3pt]
         {\scriptsize \textbf{Resources:}\;
          Neo4j/PrimeKG (20k+ diseases, 30k+ proteins)\; $\cdot$\;
          GRetriever micro-service (port~8001)\; $\cdot$\;
          OpenAI embeddings\; $\cdot$\;
          PyTorch Geometric}\\[3pt]
         {\scriptsize \textit{See companion illustration for internal
          GRetriever architecture.}}};

    \draw[flowline] (prompt.south) -- (tool.north);

    % === WRITES ===
    \node[rwpanel, draw=blue!30!black, fill=stateBlue,
          below=0.4cm of tool] (writes)
        {\textbf{$\blacktriangle$ Writes to AgentState}\\[5pt]
         $\bullet$\; \texttt{agent\_outputs[]} {\scriptsize (str)} ---
          structured KG results per gene, including:\\[2pt]
         \hphantom{$\bullet$\;}
          {\scriptsize --- pathway membership chains
          (e.g.\ \textit{CDK5 $\rightarrow$ cell-cycle signalling
          $\rightarrow$ neuronal death})}\\[2pt]
         \hphantom{$\bullet$\;}
          {\scriptsize --- GO biological-process terms}\\[2pt]
         \hphantom{$\bullet$\;}
          {\scriptsize --- drug--gene interactions and existing
          therapeutic targets}\\[2pt]
         \hphantom{$\bullet$\;}
          {\scriptsize --- protein--protein interaction partners}};

    \draw[flowline] (tool.south) -- (writes.north);

    % === DOWNSTREAM CONSUMERS ===
    \node[consumerlabel, below=0.3cm of writes] (consumers)
        {$\downarrow$\; Consumed by:\;
         \textbf{PubMedResearcher} (pathway context)\; $\cdot$\;
         \textbf{ScientistsAgent} (KG evidence for hypothesis
          anchoring)\; $\cdot$\;
         \textbf{ReportingNode} (KG analysis section)};

\end{tikzpicture}%
}% end resizebox
\caption{\textbf{BioMarkerKGAgent} --- workflow interface view.
In the typical realisation shown here this agent runs second, but
the LLM planner can reorder, skip, or re-invoke any agent at any
point---for instance, the planner may call it again after literature
search reveals new gene targets.  When invoked, it reads the
\texttt{top\_genes} list from \texttt{AgentState} and queries the
PrimeKG knowledge graph for each gene's neighbourhood, targeting the
top 10--15 DEGs to build explicit
Gene$\rightarrow$Pathway$\rightarrow$Phenotype chains.  Internally,
the tool communicates with a GRetriever micro-service that extracts
PCST subgraphs, encodes them with a Graph Attention Network, and
decodes answers with Llama-3.1-8B (see companion illustration for
architecture details).  Its output---appended to
\texttt{agent\_outputs[]}---provides the pathway and interaction
context that downstream agents use to anchor hypotheses in structured
biomedical knowledge.}
\label{fig:closeup-biomarkerkg}
\end{figure}

% ============================================================
\begin{figure}[htbp]
\centering
\resizebox{\textwidth}{!}{%
\begin{tikzpicture}[node distance=0.35cm, font=\sffamily,
    stepbar/.style={rectangle, rounded corners=6pt, thick,
        minimum width=19cm, align=center, font=\sffamily,
        drop shadow={opacity=0.10}, inner sep=6pt, text width=18cm},
    rwpanel/.style={rectangle, rounded corners=5pt, thick,
        minimum width=19cm, align=left, inner sep=6pt,
        font=\sffamily\small, text width=17.5cm},
    infopanel/.style={rectangle, rounded corners=5pt, draw=gray!50, thick,
        minimum width=19cm, align=left, fill=white,
        inner sep=6pt, font=\sffamily\small, text width=17.5cm},
    toolpanel/.style={rectangle, rounded corners=5pt, draw=gray!60, thick,
        minimum width=19cm, align=left, fill=gray!4,
        inner sep=6pt, font=\sffamily\small, text width=17.5cm},
    flowline/.style={->, >=Stealth, very thick, draw=gray!40},
    consumerlabel/.style={font=\sffamily\small\itshape, text=gray!60!black,
        align=center, text width=19cm}
]
    % === STEP BAR ===
    \node[stepbar, draw=cyan!50!black, fill=agentCyan] (step)
        {\iconGraph{cyan!60!black}\;\;
         \textbf{\large BioMarkerKGAgent --- Knowledge-Graph
          Contextualisation}\\[3pt]
         {\small Planner: \textit{``Query top 10--15 DEGs for KG
          neighbours --- builds Gene$\rightarrow$Pathway$\rightarrow$Phenotype
          chains.''}}\\[2pt]
         {\scriptsize\sffamily\color{gray!60!black}
          Shown as a second step in this example;
          the planner may invoke it earlier, later, or
          re-invoke it after replanning.}};

    % === READS ===
    \node[rwpanel, draw=green!40!black, fill=agentGreen,
          below=0.4cm of step] (reads)
        {\textbf{$\blacktriangledown$ Reads from AgentState}\\[5pt]
         $\bullet$\; \texttt{query} {\scriptsize (str)} ---
          original research question\\[2pt]
         $\bullet$\; \texttt{top\_genes} {\scriptsize (list[str])} ---
          ranked DEGs \emph{from OmicMiningAgent}\\[2pt]
         $\bullet$\; \texttt{previous\_results} {\scriptsize (str)} ---
          accumulated text output from Step~1};

    \draw[flowline] (step.south) -- (reads.north);

    % === PROMPT HIGHLIGHTS ===
    \node[infopanel, below=0.4cm of reads] (prompt)
        {\textbf{System Prompt Highlights}
         {\scriptsize (from \texttt{BioMarkerKGAgent.system\_message})}\\[5pt]
         \textbf{1.}\; ``Query the PrimeKG knowledge graph to retrieve
          relevant biomedical information''\\[2pt]
         \textbf{2.}\; ``Provide structured data and relationships from
          the knowledge graph to enhance understanding''\\[2pt]
         \textbf{3.}\; {\itshape (From planning prompt)}\;
          Query for \emph{both} up-regulated and down-regulated gene
          sets\\[2pt]
         \textbf{4.}\; {\itshape (From planning prompt)}\;
          Be THOROUGH --- query KG for multiple gene sets to build
          Gene$\rightarrow$Pathway$\rightarrow$Phenotype chains};

    \draw[flowline] (reads.south) -- (prompt.north);

    % === TOOL ===
    \node[toolpanel, below=0.4cm of prompt] (tool)
        {\textbf{Tool:}\;
         \texttt{biomarker\_kg\_tool(query)}\\[5pt]
         {\scriptsize \textbf{Internal pipeline:}\;
          text embedding $\rightarrow$
          PrimeKG subgraph retrieval (PCST) $\rightarrow$
          GAT graph encoder $\rightarrow$
          Llama-3.1-8B decoder $\rightarrow$
          ranked hypotheses}\\[3pt]
         {\scriptsize \textbf{Returns:}\;
          drug interactions, pathway memberships, GO terms,
          protein--protein interactors, disease associations}\\[3pt]
         {\scriptsize \textbf{Resources:}\;
          Neo4j/PrimeKG (20k+ diseases, 30k+ proteins)\; $\cdot$\;
          GRetriever micro-service (port~8001)\; $\cdot$\;
          OpenAI embeddings\; $\cdot$\;
          PyTorch Geometric}\\[3pt]
         {\scriptsize \textit{See companion illustration for internal
          GRetriever architecture.}}};

    \draw[flowline] (prompt.south) -- (tool.north);

    % === WRITES ===
    \node[rwpanel, draw=blue!30!black, fill=stateBlue,
          below=0.4cm of tool] (writes)
        {\textbf{$\blacktriangle$ Writes to AgentState}\\[5pt]
         $\bullet$\; \texttt{agent\_outputs[]} {\scriptsize (str)} ---
          structured KG results per gene, including:\\[2pt]
         \hphantom{$\bullet$\;}
          {\scriptsize --- pathway membership chains
          (e.g.\ \textit{CDK5 $\rightarrow$ cell-cycle signalling
          $\rightarrow$ neuronal death})}\\[2pt]
         \hphantom{$\bullet$\;}
          {\scriptsize --- GO biological-process terms}\\[2pt]
         \hphantom{$\bullet$\;}
          {\scriptsize --- drug--gene interactions and existing
          therapeutic targets}\\[2pt]
         \hphantom{$\bullet$\;}
          {\scriptsize --- protein--protein interaction partners}};

    \draw[flowline] (tool.south) -- (writes.north);

    % === DOWNSTREAM CONSUMERS ===
    \node[consumerlabel, below=0.3cm of writes] (consumers)
        {$\downarrow$\; Consumed by:\;
         \textbf{PubMedResearcher} (pathway context)\; $\cdot$\;
         \textbf{ScientistsAgent} (KG evidence for hypothesis
          anchoring)\; $\cdot$\;
         \textbf{ReportingNode} (KG analysis section)};

\end{tikzpicture}%
}% end resizebox
\caption{\textbf{BioMarkerKGAgent} --- workflow interface view.
In the typical realisation shown here this agent runs second, but
the LLM planner can reorder, skip, or re-invoke any agent at any
point---for instance, the planner may call it again after literature
search reveals new gene targets.  When invoked, it reads the
\texttt{top\_genes} list from \texttt{AgentState} and queries the
PrimeKG knowledge graph for each gene's neighbourhood, targeting the
top 10--15 DEGs to build explicit
Gene$\rightarrow$Pathway$\rightarrow$Phenotype chains.  Internally,
the tool communicates with a GRetriever micro-service that extracts
PCST subgraphs, encodes them with a Graph Attention Network, and
decodes answers with Llama-3.1-8B (see companion illustration for
architecture details).  Its output---appended to
\texttt{agent\_outputs[]}---provides the pathway and interaction
context that downstream agents use to anchor hypotheses in structured
biomedical knowledge.}
\label{fig:closeup-biomarkerkg}
\end{figure}

% ============================================================
\begin{figure}[htbp]
\centering
\resizebox{\textwidth}{!}{%
\begin{tikzpicture}[node distance=0.35cm, font=\sffamily,
    stepbar/.style={rectangle, rounded corners=6pt, thick,
        minimum width=19cm, align=center, font=\sffamily,
        drop shadow={opacity=0.10}, inner sep=6pt, text width=18cm},
    rwpanel/.style={rectangle, rounded corners=5pt, thick,
        minimum width=19cm, align=left, inner sep=6pt,
        font=\sffamily\small, text width=17.5cm},
    infopanel/.style={rectangle, rounded corners=5pt, draw=gray!50, thick,
        minimum width=19cm, align=left, fill=white,
        inner sep=6pt, font=\sffamily\small, text width=17.5cm},
    toolpanel/.style={rectangle, rounded corners=5pt, draw=gray!60, thick,
        minimum width=19cm, align=left, fill=gray!4,
        inner sep=6pt, font=\sffamily\small, text width=17.5cm},
    flowline/.style={->, >=Stealth, very thick, draw=gray!40},
    consumerlabel/.style={font=\sffamily\small\itshape, text=gray!60!black,
        align=center, text width=19cm}
]
    % === STEP BAR ===
    \node[stepbar, draw=cyan!50!black, fill=agentCyan] (step)
        {\iconGraph{cyan!60!black}\;\;
         \textbf{\large BioMarkerKGAgent --- Knowledge-Graph
          Contextualisation}\\[3pt]
         {\small Planner: \textit{``Query top 10--15 DEGs for KG
          neighbours --- builds Gene$\rightarrow$Pathway$\rightarrow$Phenotype
          chains.''}}\\[2pt]
         {\scriptsize\sffamily\color{gray!60!black}
          Shown as a second step in this example;
          the planner may invoke it earlier, later, or
          re-invoke it after replanning.}};

    % === READS ===
    \node[rwpanel, draw=green!40!black, fill=agentGreen,
          below=0.4cm of step] (reads)
        {\textbf{$\blacktriangledown$ Reads from AgentState}\\[5pt]
         $\bullet$\; \texttt{query} {\scriptsize (str)} ---
          original research question\\[2pt]
         $\bullet$\; \texttt{top\_genes} {\scriptsize (list[str])} ---
          ranked DEGs \emph{from OmicMiningAgent}\\[2pt]
         $\bullet$\; \texttt{previous\_results} {\scriptsize (str)} ---
          accumulated text output from Step~1};

    \draw[flowline] (step.south) -- (reads.north);

    % === PROMPT HIGHLIGHTS ===
    \node[infopanel, below=0.4cm of reads] (prompt)
        {\textbf{System Prompt Highlights}
         {\scriptsize (from \texttt{BioMarkerKGAgent.system\_message})}\\[5pt]
         \textbf{1.}\; ``Query the PrimeKG knowledge graph to retrieve
          relevant biomedical information''\\[2pt]
         \textbf{2.}\; ``Provide structured data and relationships from
          the knowledge graph to enhance understanding''\\[2pt]
         \textbf{3.}\; {\itshape (From planning prompt)}\;
          Query for \emph{both} up-regulated and down-regulated gene
          sets\\[2pt]
         \textbf{4.}\; {\itshape (From planning prompt)}\;
          Be THOROUGH --- query KG for multiple gene sets to build
          Gene$\rightarrow$Pathway$\rightarrow$Phenotype chains};

    \draw[flowline] (reads.south) -- (prompt.north);

    % === TOOL ===
    \node[toolpanel, below=0.4cm of prompt] (tool)
        {\textbf{Tool:}\;
         \texttt{biomarker\_kg\_tool(query)}\\[5pt]
         {\scriptsize \textbf{Internal pipeline:}\;
          text embedding $\rightarrow$
          PrimeKG subgraph retrieval (PCST) $\rightarrow$
          GAT graph encoder $\rightarrow$
          Llama-3.1-8B decoder $\rightarrow$
          ranked hypotheses}\\[3pt]
         {\scriptsize \textbf{Returns:}\;
          drug interactions, pathway memberships, GO terms,
          protein--protein interactors, disease associations}\\[3pt]
         {\scriptsize \textbf{Resources:}\;
          Neo4j/PrimeKG (20k+ diseases, 30k+ proteins)\; $\cdot$\;
          GRetriever micro-service (port~8001)\; $\cdot$\;
          OpenAI embeddings\; $\cdot$\;
          PyTorch Geometric}\\[3pt]
         {\scriptsize \textit{See companion illustration for internal
          GRetriever architecture.}}};

    \draw[flowline] (prompt.south) -- (tool.north);

    % === WRITES ===
    \node[rwpanel, draw=blue!30!black, fill=stateBlue,
          below=0.4cm of tool] (writes)
        {\textbf{$\blacktriangle$ Writes to AgentState}\\[5pt]
         $\bullet$\; \texttt{agent\_outputs[]} {\scriptsize (str)} ---
          structured KG results per gene, including:\\[2pt]
         \hphantom{$\bullet$\;}
          {\scriptsize --- pathway membership chains
          (e.g.\ \textit{CDK5 $\rightarrow$ cell-cycle signalling
          $\rightarrow$ neuronal death})}\\[2pt]
         \hphantom{$\bullet$\;}
          {\scriptsize --- GO biological-process terms}\\[2pt]
         \hphantom{$\bullet$\;}
          {\scriptsize --- drug--gene interactions and existing
          therapeutic targets}\\[2pt]
         \hphantom{$\bullet$\;}
          {\scriptsize --- protein--protein interaction partners}};

    \draw[flowline] (tool.south) -- (writes.north);

    % === DOWNSTREAM CONSUMERS ===
    \node[consumerlabel, below=0.3cm of writes] (consumers)
        {$\downarrow$\; Consumed by:\;
         \textbf{PubMedResearcher} (pathway context)\; $\cdot$\;
         \textbf{ScientistsAgent} (KG evidence for hypothesis
          anchoring)\; $\cdot$\;
         \textbf{ReportingNode} (KG analysis section)};

\end{tikzpicture}%
}% end resizebox
\caption{\textbf{BioMarkerKGAgent} --- workflow interface view.
In the typical realisation shown here this agent runs second, but
the LLM planner can reorder, skip, or re-invoke any agent at any
point---for instance, the planner may call it again after literature
search reveals new gene targets.  When invoked, it reads the
\texttt{top\_genes} list from \texttt{AgentState} and queries the
PrimeKG knowledge graph for each gene's neighbourhood, targeting the
top 10--15 DEGs to build explicit
Gene$\rightarrow$Pathway$\rightarrow$Phenotype chains.  Internally,
the tool communicates with a GRetriever micro-service that extracts
PCST subgraphs, encodes them with a Graph Attention Network, and
decodes answers with Llama-3.1-8B (see companion illustration for
architecture details).  Its output---appended to
\texttt{agent\_outputs[]}---provides the pathway and interaction
context that downstream agents use to anchor hypotheses in structured
biomedical knowledge.}
\label{fig:closeup-biomarkerkg}
\end{figure}
